% !PRINTED 212
\weberpage{212}
%<<CONT
\weq{10}{S[F_0(x)] = 0, \quad S[F_1(x)] = 0, \ \ldots \ S[F_{n-2}(x)] = 0.}
Setzen wir nach (3)
\begin{align*}
F(t, x) &= \frac{f(t)}{t-x} - \frac{1}{n} f'(t) \tag{11}\\
&= t^{n-2} F_0(x) + t^{n-3} F_1(x) + \cdots + t F_{n-3}(x) + F_{n-2}(x),
\end{align*}
so geht die rechte Seite von (8) aus $F(t, x)$ hervor durch die
Ersetzung von $t^k$ durch $t_k$.

Nehmen wir von vornherein $y$ in der Form
\weq{12}{y = t_{n-2} F_0(x) + t_{n-3} F_1(x) + \cdots + t_0 F_{n-2}(x)}
an, so ist die Gleichung
\weq{13}{S(y) = 0}
identisch befriedigt. Welchen Nutzen diese Form der Substitution
gew\"ahrt, werden die n\"achsten Betrachtungen zeigen.

\begin{center}
\S.\,69.

\gesperrt{Invarianten-Eigenschaft der Tschirnhausen-Transformation.}
\end{center}

Es ist jetzt der Einfluss zu untersuchen, den eine lineare
Transformation, der wir die Function $f(x)$ unterwerfen, auf die
Tschirnhausen-Transformation hat.

Wir machen in $f(x)$ die lineare Substitution
\weq{1}{x = \frac{\alpha \xi + \beta}{\gamma \xi + \delta}, \quad \alpha \delta - \beta \gamma = r,}
wodurch wir erhalten
\weq{2}{\varphi(\xi) = (\gamma \xi + \delta)^n f\left(\frac{\alpha \xi + \beta}{\gamma \xi + \delta}\right),}
so dass $\varphi(\xi)$ eine ganze rationale Function $n^{\text{ten}}$ Grades und
$\varphi(\xi) = 0$ die durch (1) transformirte Gleichung $f(x) = 0$ ist.

Wir leiten nun eine Function $\Phi(\tau, \xi)$ ganz in derselben
Weise aus $\varphi(\xi)$ ab, wie wir im vorigen Paragraphen $F(t, x)$ aus
$f(x)$ abgeleitet haben, n\"amlich
\begin{align*}
\Phi(\tau, \xi) &= \frac{\varphi(\tau)}{\tau - \xi} - \frac{1}{n} \varphi'(\tau) \tag{3}\\
&= \tau^{n-2} \Phi_0(\xi) + \tau^{n-3} \Phi_1(\xi) + \cdots + \tau \Phi_{n-3}(\xi) + \Phi_{n-2}(\xi),
\end{align*}
worin $\tau$ ebenso wie $t$ eine Variable ist und die Functionen
$\Phi_0(\xi), \Phi_1(\xi) \ldots \Phi_{n-2}(\xi)$, ebenso aus $\xi$ und den Co\"efficienten
%CONT>>
